% GENERATED FILE. Edit canonical lessons, then run npm run generate:books.
% canonical-source-hash: fnv1a64:8bc82509088f86a9
% canonical-lessons: MR-C495-putanya, MR-C495-bhaca, MR-C495-natu, MR-C495-panjoba, MR-C495-atya

\chapter{Nephews and Grandchildren}
\label{ch:mr-nephews-and-grandchildren}
\hlchaptermodality{495}

\begin{chapteropening}
\textbf{By the end of this chapter:} \emph{I can name a nephew, a grandson, a great-grandfather and a father's sister.}

Complete the last lesson of chapter 495: I can name a nephew, a grandson, a great-grandfather and a father's sister.
\end{chapteropening}

\section[\texorpdfstring{\mr{पुतण्या} (putaṇyā)}{पुतण्या (putaṇyā)}]{\mr{पुतण्या} (putaṇyā) — a brother's son}

\label{lesson:MR-C495-putanya}



\begin{quote}
Before the new one: say the Marathi for an orphan, then the Marathi for full-blood.
\end{quote}

\subsection*{You'll want to know: \mr{पुतण्या}}
\textbf{\mr{पुतण्या}} — \emph{putaṇyā} — \textquotedblleft{}a brother's son\textquotedblright{}.

\begin{quote}
\mr{माझा} \mr{पुतण्या} \mr{पुण्यात} \mr{शिकतो}.
\end{quote}

\begin{quote}
\emph{mājhā putaṇyā puṇyāt śikto.}
\end{quote}

\begin{quote}
My nephew studies in Pune.
\end{quote}

Masculine.

A brother's son; a brother's daughter is \mr{पुतणी}.

\begin{cousinweb}
\textbf{\mr{पुतण्या}} (\emph{putaṇyā}) is inherited: CDIAL lists Marathi \emph{putṇyā}, \textquotedblleft{}brother's son\textquotedblright{}, under the reconstructed \emph{*pitriyajāni}, \textquotedblleft{}wife of one of the father's family\textquotedblright{} (CDIAL 8189).
\end{cousinweb}

\begin{grammarlens}[title={What you've built}]
One more word for describing what you lived and felt.
\end{grammarlens}

\subsection*{Your turn}
\begin{itemize}
\raggedright
  \item \emph{Say it:} \emph{putaṇyā}
  \item \emph{Say it:} \emph{putaṇyā}, once more
  \item \emph{Say it:} \emph{mājhā putaṇyā puṇyāt śikto.} — My nephew studies in Pune.
  \item \emph{Recall:} say the Marathi for an orphan, then the Marathi for full-blood, then say \emph{putaṇyā} again
\end{itemize}

\begin{tcolorbox}[breakable,colback=teal!4,colframe=teal!35!black,title={Before you move on}]
What does \mr{पुतण्या} mean? (\textquotedblleft{}A brother's son\textquotedblright{}.) Say it once more.
\end{tcolorbox}
\section[\texorpdfstring{\mr{भाचा} (bhācā)}{भाचा (bhācā)}]{\mr{भाचा} (bhācā) — a sister's son}

\label{lesson:MR-C495-bhaca}



\begin{quote}
Before the new one: say the Marathi for full-blood, then the Marathi for a brother's son.
\end{quote}

\subsection*{You'll want to know: \mr{भाचा}}
\textbf{\mr{भाचा}} — \emph{bhācā} — \textquotedblleft{}a sister's son\textquotedblright{}.

\begin{quote}
\mr{सुट्टीत} \mr{माझा} \mr{भाचा} \mr{आमच्याकडे} \mr{राहायला} \mr{आला}.
\end{quote}

\begin{quote}
\emph{suṭṭīt mājhā bhācā āmcyākaḍe rāhāylā ālā.}
\end{quote}

\begin{quote}
My nephew came to stay with us in the holidays.
\end{quote}

Masculine.

A sister's son; \mr{भाची}, a sister's daughter. To the child, the mother's brother is \mr{मामा}.

\begin{cousinweb}
\textbf{\mr{भाचा}} (\emph{bhācā}) is filed by CDIAL under \emph{bhrātrya}, \textquotedblleft{}pertaining to a brother\textquotedblright{}, together with Prakrit \emph{bhacca}, \textquotedblleft{}sister's son\textquotedblright{} (CDIAL 9673); how a brother's word came to name a sister's son is not settled.
\end{cousinweb}

\begin{grammarlens}[title={What you've built}]
One more word for describing what you lived and felt.
\end{grammarlens}

\subsection*{Your turn}
\begin{itemize}
\raggedright
  \item \emph{Say it:} \emph{bhācā}
  \item \emph{Say it:} \emph{bhācā}, once more
  \item \emph{Say it:} \emph{suṭṭīt mājhā bhācā āmcyākaḍe rāhāylā ālā.} — My nephew came to stay with us in the holidays.
  \item \emph{Recall:} say the Marathi for full-blood, then the Marathi for a brother's son, then say \emph{bhācā} again
\end{itemize}

\begin{tcolorbox}[breakable,colback=teal!4,colframe=teal!35!black,title={Before you move on}]
What does \mr{भाचा} mean? (\textquotedblleft{}A sister's son\textquotedblright{}.) Say it once more.
\end{tcolorbox}
\section[\texorpdfstring{\mr{नातू} (nātū)}{नातू (nātū)}]{\mr{नातू} (nātū) — a grandson}

\label{lesson:MR-C495-natu}



\begin{quote}
Before the new one: say the Marathi for a brother's son, then the Marathi for a sister's son.
\end{quote}

\subsection*{You'll want to know: \mr{नातू}}
\textbf{\mr{नातू}} — \emph{nātū} — \textquotedblleft{}a grandson\textquotedblright{}.

\begin{quote}
\mr{आजोबा} \mr{नातवाबरोबर} \mr{बागेत} \mr{फिरतात}.
\end{quote}

\begin{quote}
\emph{ājobā nātvābarobar bāget phirtāt.}
\end{quote}

\begin{quote}
Grandfather walks in the garden with his grandson.
\end{quote}

Masculine.

Oblique \mr{नातव}-: \mr{नातवाला}. \mr{नात}, a granddaughter.

\begin{cousinweb}
\textbf{\mr{नातू}} (\emph{nātū}) is a tadbhava, inherited from Sanskrit \emph{naptṛ}, \textquotedblleft{}grandson\textquotedblright{}; CDIAL lists Marathi \emph{nātū} (CDIAL 6955b).
\end{cousinweb}

\begin{grammarlens}[title={What you've built}]
One more word for describing what you lived and felt.
\end{grammarlens}

\subsection*{Your turn}
\begin{itemize}
\raggedright
  \item \emph{Say it:} \emph{nātū}
  \item \emph{Say it:} \emph{nātū}, once more
  \item \emph{Say it:} \emph{ājobā nātvābarobar bāget phirtāt.} — Grandfather walks in the garden with his grandson.
  \item \emph{Recall:} say the Marathi for a brother's son, then the Marathi for a sister's son, then say \emph{nātū} again
\end{itemize}

\begin{tcolorbox}[breakable,colback=teal!4,colframe=teal!35!black,title={Before you move on}]
What does \mr{नातू} mean? (\textquotedblleft{}A grandson\textquotedblright{}.) Say it once more.
\end{tcolorbox}
\section[\texorpdfstring{\mr{पणजोबा} (paṇjobā)}{पणजोबा (paṇjobā)}]{\mr{पणजोबा} (paṇjobā) — a great-grandfather}

\label{lesson:MR-C495-panjoba}



\begin{quote}
Before the new one: say the Marathi for a sister's son, then the Marathi for a grandson.
\end{quote}

\subsection*{You'll want to know: \mr{पणजोबा}}
\textbf{\mr{पणजोबा}} — \emph{paṇjobā} — \textquotedblleft{}a great-grandfather\textquotedblright{}.

\begin{quote}
\mr{माझे} \mr{पणजोबा} \mr{शंभर} \mr{वर्षं} \mr{जगले}.
\end{quote}

\begin{quote}
\emph{mājhe paṇjobā śambhar varṣaṁ jagle.}
\end{quote}

\begin{quote}
My great-grandfather lived a hundred years.
\end{quote}

Masculine.

Built like \mr{आजोबा}; \mr{पणजी}, a great-grandmother.

\begin{cousinweb}
\textbf{\mr{पणजोबा}} (\emph{paṇjobā}) is a Marathi kin word with the prefix \mr{पण}-, "great-", the one CDIAL shows in Marathi \emph{paṇtū}, \textquotedblleft{}great-grandson\textquotedblright{}, under \emph{*praṇaptṛ} (CDIAL 8523); the rest of the word is usually traced to the family of \mr{आजोबा}.
\end{cousinweb}

\begin{grammarlens}[title={What you've built}]
One more word for describing what you lived and felt.
\end{grammarlens}

\subsection*{Your turn}
\begin{itemize}
\raggedright
  \item \emph{Say it:} \emph{paṇjobā}
  \item \emph{Say it:} \emph{paṇjobā}, once more
  \item \emph{Say it:} \emph{mājhe paṇjobā śambhar varṣaṁ jagle.} — My great-grandfather lived a hundred years.
  \item \emph{Recall:} say the Marathi for a sister's son, then the Marathi for a grandson, then say \emph{paṇjobā} again
\end{itemize}

\begin{tcolorbox}[breakable,colback=teal!4,colframe=teal!35!black,title={Before you move on}]
What does \mr{पणजोबा} mean? (\textquotedblleft{}A great-grandfather\textquotedblright{}.) Say it once more.
\end{tcolorbox}
\section[\texorpdfstring{\mr{आत्या} (ātyā)}{आत्या (ātyā)}]{\mr{आत्या} (ātyā) — a father's sister}

\label{lesson:MR-C495-atya}



\begin{quote}
Before the new one: say the Marathi for a grandson, then the Marathi for a great-grandfather.
\end{quote}

\subsection*{You'll want to know: \mr{आत्या}}
\textbf{\mr{आत्या}} — \emph{ātyā} — \textquotedblleft{}a father's sister\textquotedblright{}.

\begin{quote}
\mr{माझी} \mr{आत्या} \mr{गावात} \mr{राहते}.
\end{quote}

\begin{quote}
\emph{mājhī ātyā gāvāt rāhte.}
\end{quote}

\begin{quote}
My father's sister lives in the village.
\end{quote}

Feminine.

A father's sister; a mother's sister is \mr{मावशी}.

\begin{cousinweb}
\textbf{\mr{आत्या}} (\emph{ātyā}) is inherited: CDIAL lists Marathi \emph{āt}, \emph{ātū} and \emph{āte}, \textquotedblleft{}father's sister\textquotedblright{}, under Sanskrit \emph{attā}, \textquotedblleft{}mother, elder sister\textquotedblright{} (CDIAL 222), a word usually traced to Dravidian, as CDIAL does: Tamil \emph{attai}, \textquotedblleft{}aunt\textquotedblright{}.
\end{cousinweb}

\begin{grammarlens}[title={What you've built}]
One more word for describing what you lived and felt.
\end{grammarlens}

\subsection*{Your turn}
\begin{itemize}
\raggedright
  \item \emph{Say it:} \emph{ātyā}
  \item \emph{Say it:} \emph{ātyā}, once more
  \item \emph{Say it:} \emph{mājhī ātyā gāvāt rāhte.} — My father's sister lives in the village.
  \item \emph{Recall:} say the Marathi for a grandson, then the Marathi for a great-grandfather, then say \emph{ātyā} again
\end{itemize}

\begin{tcolorbox}[breakable,colback=teal!4,colframe=teal!35!black,title={Before you move on}]
What does \mr{आत्या} mean? (\textquotedblleft{}A father's sister\textquotedblright{}.) Say it once more.
\end{tcolorbox}
